\chapter*{Introduction}
\addcontentsline{toc}{chapter}{Introduction}
%
\noindent
%
\textcolor{red}{ESTO NO LO HE REVISADO AUN}
\noindent The objective of understanding the fundamental laws that govern matter and its interactions has guided the evolution of particle physics from its earliest concepts to its most sophisticated theories. Discovery of novel phenomena gave rise to the development of theoretical frameworks that emerged to describe the behavior of elementary particles and forces. The theoretical synthesis achieved in the last century culminated in the Standard Model of particle physics (SM), a quantum field theory that elegantly describes the structure of all known particles and unifies the electromagnetic, weak and strong interactions. 

The SM theory has been met with notable success due to its predictive accuracy and the extensive supporting experimental evidence, thanks to collider experiments such as the Large Hadron Collider (LHC). The final piece in the experimental exploration of the SM model, the Higgs boson, was proposed in 1964 and discovered in 2012 by the CMS and ATLAS collaborations [ref], confirming the electroweak symmetry breaking (EWSB) through the Brout-Englert-Higgs (BEH) mechanism. The discovery of the Higgs boson marked an exciting milestone in particle physics, not only because it completes the Standard Model, making it a theory valid up to the Plank scale, but because of its unique properties.

As the only known elementary spin-0 boson, its observation opened a new frontier of exploration into the scalar sector, and ushered the LHC physics program into a new era. During the LHC Run 1 (2009-2013), which allowed for the discovery of the Higgs boson, the first measurements of its properties were conducted and found to be consistent with SM predictions. After Run I completion, a new phase of characterization and precision measurements focused on the mass and with of the Higgs boson, and its couplings to other SM particles began during Run 2 (2015-2018). The exploration of the Higgs sector across all probed features has shown that the Higgs particle behaves in a very SM-like manner thus far.

However, the overall picture is incomplete. Although the Standard Model validly describes all the measurements performed at the GeV and TeV energy scales, several fundamental questions remain still unanswered. Quests as the matter-antimatter asymmetry, the nature of dark matter and gravity, among others, indicate the limitations of the SM, and suggest that it is only a part of a more comprehensive physics model. The existence of a so-called Beyond the Standard Model (BSM) physics is the enigma that has become the current focus of both theoretical and experimental efforts. Many models have been proposed to address these open issues. In this context, exploring the Higgs sector becomes a particularly promising window into new physics given the unresolved questions surrounding the stability of the electroweak vacuum and the naturalness of the Higgs mass. 

One of the key elements remaining in the exploration of the SM Higgs sector is the Higgs boson self-interaction, encoded in the trilinear self-coupling $\lambda_{HHH}$, which has not yet been characterized by the data collected to date by the LHC experiments to date. The crucial role of the Higgs self-interaction lies in its intrinsic relation to the shape of the Higgs potential and its non-zero vacuum expectation value (VEV). Thus, measuring it would be a powerful closure test for the EWSB and the SM. The most direct way to access the Higgs self-interaction is through the Higgs boson pair (HH) production. As a direct link to the scalar potential, HH production is highly sensitive to BSM effects, which manifest either as quantum loops contributions that induce deviations in the production rate and kinematics, or via the resonant decay of new states into a Higgs boson pair. Because the HH process explores the fundamental nature of the Higgs scalar field as it provides unique handle for directly accessing $\lambda_{HHH}$, its measurement has become a fundamental priority for the particle physics community. 
    
And yet, as previously mentioned, the experimental measurement of the trilinear self-coupling remains an outstanding task: observing HH SM production at the LHC is extremely challenging due to its small cross section. However, the amount of data collected during the operation of the LHC does provide a measure of the sensitivity achieved for this process. Through the study and combination of multiple Higgs boson decay modes, stringent upper limits on the HH production cross section and exclusion ranges for the self-coupling have already been derived using LHC Run 2 data. In this regard, the increase in statistics provided by the LHC Run 3 (2022-2026) will entail substantial benefits for investigating the Di-Higgs mechanism. Among the various decay channels of interest, the $bb\tau\tau$ final state stands out thanks to its relatively large branching fraction and the relatively low background contamination, making it one of the most sensitive channels for HH searches. 

Despite the statistical limitations that are gradually being overcome thanks to improved reconstruction techniques and identification methods, the definitive breakthrough will come from the High-Luminosity LHC (HL-LHC). Scheduled to begin operation in 2029, the HL-LHC is expected to significantly increase in instantaneous luminosity and collected data, paving the way for previously unexplored areas of physics. In particular, in the scalar sector, it is expected to provide the first evidence of Di-Higgs production, allowing for a deeper exploration of the EWSB properties and the SM validity.

The work presented in this thesis contributes to the broader effort to probe the Higgs sector by studying the $bb\tau\tau$ decay channel, aiming to overcome the experimental challenges posed by this line of research. Along with the limited statistics handicap, the sensitivity is also highly constrained by correctly identifying the objects involved. The high rate of proton bunches crossing in the center of the detector creates an extremely dense environment in which millions of collisions occur every millisecond. Proper identification of the objects of interest and rejection of the abundant jet background are therefore crucial to achieving maximal sensitivity. Thus, this thesis particularly emphasizes the identification of the $b$ quarks and $\tau$ leptons involved in the decay. 


The thesis is organized in three parts. The first one sets the framework of the work. Chapter \ref{ch:theory} introduces the SM, the BEH mechanism and the phenomenology of Higgs boson pair production, together with the experimental status of its searches, and Chapter \ref{ch:LHC_CMS} describes the LHC, the CMS detector and the reconstruction of the physics objects.

The second part presents the search for Higgs boson pair production in the $bb\tau\tau$ final state with the Run 3 data, preceded by the identification of the objects to which this thesis has contributed most directly. Chapter \ref{ch:Tau} describes the identification of hadronically decaying $\tau$ leptons and its validation with the Run 3 data, which constitutes the contribution of this thesis to the Tau POG and has been published in Ref.~\cite{CMS-DP-2025-074}. Chapter \ref{ch:taggers} presents the algorithms that assign the jets of the event to the $H\to b\bar{b}$ candidate and to the VBF topology. The retraining of \textsc{HH-BTag} for the Run 2 search for resonant $HH$ production \cite{CMS-PAS-B2G-24-011}, its validation for the study of $ZZ$ and $ZH$ production in the same final state \cite{CMS:2026zzzh}, its retraining for the Run 3 analysis and the development of \textsc{VBFjTag} are contributions of this thesis. Chapters \ref{ch:bbtt_analysis} and \ref{ch:bbtt_results} describe the Run 3 analysis \cite{CMS-PAS-HIG-25-008} and its results, including their combination with the Run 2 dataset, to which this thesis has contributed \textcolor{red}{[completar: \textsc{HH-BTag} y \textsc{VBFjTag}, SFs QuadJet, validación del método ABCD para el Run2 donde lo meto? ...]}.

The third part addresses the interpretation of $HH$ searches beyond the SM. Chapter \ref{ch:eft_modeling} develops the modeling of the gluon fusion $HH$ signal in the HEFT framework, extending the matrix inversion method beyond the coverage of previous ATLAS and CMS results, which constitutes the main original contribution of this thesis to the BSM interpretation of $HH$ searches and is documented in Ref.~\cite{LHCHWG-2026-015}.


%Chapter \ref{1} establishes the theoretical context of this work by introducing the SM and the Di-Higgs production. The experimental context of the analysis is provided in Chapter \ref{2}, covering the LHC accelerator and the CMS experiment. Chapter \ref{3} presents the particle reconstruction and identification methods used in CMS analyses, with a focus on my contributions to $\tau$ reconstruction. The subsequent chapters concentrate on the $bb\tau\tau$ analysis: Chapter \ref{4} discusses the analysis strategy and illustrates the techniques developed to select the different objects in the final state, while the results are shown in Chapter \ref{5}

